\begin{proof}
Using the notations from Subsection \ref{subsec:inv}, we denote by $K$, the cokernel of the morphism~$\sigma^2(P_1)$,
\begin{gather} K=\frac{T^*\otimes \left(T^*_v \oplus
\Lambda^2T^*_v \oplus \Lambda^2 T^*_v\right)}{\operatorname{Im}
\sigma^2(P_1)}. \label{spacek} \end{gather}
We will prove the theorem by using the following classical result of
homological algebra, see \cite[Proposition~1.1]{grifone00}. There
exists a morphism $\varphi: R^1(P_1) \to K$ such that the sequence
\[
 R^2(P_1)\overset{\overline{\pi}_1}{\longrightarrow} R^1(P_1)
\overset{\varphi}{\longrightarrow} K
\] is exact. In particular, the
morphism  $\overline{\pi}_1$ is onto if and only if~$\varphi=0$.

We will build the morphism $\varphi$ and show that for $\theta \in
\Lambda^1_v$ such that $j^1_u\theta \in R^1_u(P_1)\subset J^1_uT^*_v$, a~f\/irst-order
solution of~$P_1$ at $u\in TM\setminus\{0\}$ we have that
$\varphi_u\theta=0$ if and only if~$(d_R\theta)_u=0$. The morphism
$\varphi$ is represented in the diagram~\eqref{diagram} by dashed
arrows.

To build $\varphi$, we have to def\/ine f\/irst a morphism of vector bundles
\[
 \tau: \ T^*\otimes \left(T^*_v \oplus
\Lambda^2T^*_v \oplus \Lambda^2 T^*_v\right) \to K,
\] such that the
f\/irst row in the following diagram is exact.
\begin{gather}
  \label{diagram}
  {\diagram
 &  0 \dto  & 0 \dto & 0 \dto
  \\
 0\rto & g^2(P_1) \rto \dto & S^{2}T^*\otimes T^*_v \rto^{\sigma^2(P_1)} \dto^{\varepsilon} & T^* \otimes
  F_1 \rto \rdashed<1ex>|>{\tip}^{\tau} \dto^{\varepsilon} &
  %%
  K \rto & 0
  \\
 0 \rto &   R^{2}(P_1) \rto^{i} \dto^{\overline {\pi}_{1}} & J_{2} T^*_v \rto \dto^{\pi_{1}}
  \rdashed<1ex>|>{\tip}^{p^1(P_1)}
  & J^1 F_1 \dto^{\pi} \udashed<1ex>|>{\tip}^{\nabla}
  \\
  0 \rto & R^1(P_1) \rto\rdashed<1ex>|>{\tip}^i  & J^1T^*_v \rto^{p^o(P_1)}
  \udashed<1ex>|>{\tip}  \dto & F_1 \dto
  \\
&  & 0 & 0
  \enddiagram}
\end{gather}

For the vector bundle $K$ given by formula \eqref{spacek}, the
dimension of its f\/ibres is $n^2(n-1)/2$. Therefore, we can view this
vector bundle over $TM\setminus\{0\}$ as follows
\[
K=\oplus^{(2)}\Lambda^2T^*_v \oplus^{(3)}\Lambda^3T^*_v.
\]
Therefore the map $\tau$ has $5$ components.
The f\/ive components of $\tau=(\tau_1,\dots , \tau_5)$, are given as
follows
\begin{alignat*}{4}
& \tau_{1}(A,B_1,B_2) = \tau_JA - i_{\mathbb{C}}B_1, \qquad &&
\tau_{2}(A,B_1,B_2) = \tau_hA - i_{\mathbb{C}}B_2, \qquad &&
\tau_{3}(A,B_1,B_2) = \tau_JB_1,&  \\
& \tau_{4}(A,B_1,B_2) = \tau_hB_2, \qquad &&
\tau_{5}(A,B_1,B_2) = \tau_hB_1+\tau_JB_2, &&&
\end{alignat*}
for $A\in T^*\otimes T^*_v$, $B_1, B_2 \in T^*\otimes
\Lambda^2T^*_v$. Using the above def\/inition of the f\/ive components of $\tau$,
formula \eqref{sigma2p1} that def\/ines the three components of
$\sigma^2(P_1)$, and the symmetry in the f\/irst two arguments of an
element $B\in S^2T^*\otimes T^*_v$ we can prove $\left(\tau\circ
\sigma^2(P_1)\right)(B)=0$. For example, the f\/irst component of this
composition is given by {\sloppy
\begin{gather*}
\left(\tau_1\circ \sigma^2(P_1)\right)(B)(X,Y)  =
\left(\tau_J\sigma^2({\mathcal{L}_{\mathbb{C}}})\right)(B)(X,Y)-
\left(i_{\mathbb{C}}\sigma^2(d_J)\right)(B)(X,Y)\\
 \phantom{\left(\tau_1\circ \sigma^2(P_1)\right)(B)(X,Y)}{}
 = B(JX, \mathbb{C}, Y)- B(JY, \mathbb{C}, X) - B(\mathbb{C}, JX, Y)+ B(\mathbb{C}, JY, X)=0.\!
\end{gather*} It follows that $\operatorname{Im} \sigma^2(P_1)
\subset \operatorname{Ker} \tau $. By comparing the dimensions, it
is easy to see that  \mbox{$\operatorname{Im} \sigma^2(P_1) =
\operatorname{Ker} \tau $}, and therefore the f\/irst row in diagram~\eqref{diagram} is exact.

}

Consider $\nabla$ a linear connection on $TM\setminus\{0\}$ such
that~$\nabla J=0$. It follows that the connection~$\nabla$ preserves
the vertical distribution and hence it will preserve semi-basic
forms. Therefore, one can view $\nabla$ as a connection on the
f\/ibre bundle $F_1\to TM\setminus\{0\}$.  Using Lemma~\ref{lem:nablaomega}, it follows that derivations
$\mathcal{D}_{\mathbb{C}}=i_{\mathbb{C}}\nabla$,
$\mathcal{D}_{J}=\tau_{J}\nabla$, and $\mathcal{D}_{h}=\tau_{h}\nabla$
preserve semi-basic forms. As a~f\/irst-order partial
dif\/ferential operator, we can identify connection $\nabla$ with
the bundle morphism $p^0(\nabla): J^1F_1 \to T^*\otimes F_1$.
We will use this bundle morphism to def\/ine the map $\varphi: R^1(P_1)
\to K$ we mentioned at the beginning of the proof.

Consider $\theta\in \Lambda^1_v$ such that $j^1_u\theta\in
R^1_u(P_1)\subset J_u^1T^*_v$ is a f\/irst-order solution of $P_1$
at $u\in TM\setminus\{0\}$. Then, we def\/ine
\[
 \varphi_u\theta=\tau_u\nabla P_1\theta=\tau_u(\nabla \mathcal{L}_{\mathbb{C}}\theta,
\nabla d_J\theta, \nabla d_h\theta).
\] We will compute now the
f\/ive components of map $\varphi$. Since
$\mathcal{L}_{\mathbb{C}}\theta$, $d_J\theta$, and $d_h\theta$
vanish at $u\in TM\setminus\{0\}$, using Lemma \ref{lem:dec_DL},
it follows that when acting on this semi-basic forms we have
$\mathcal{D}_{\mathbb{C}}=\mathcal{L}_{\mathbb{C}}$,
$\mathcal{D}_{J}=d_J$, and $\mathcal{D}_{h}=d_h$. Using the fact that
$[J, \mathbb{C}]=J$, $[h, \mathbb{C}]=0$,  $[J, J]=0$, and $[h,J]=0$,
it follows that
\begin{gather*}
\tau_{1}\left(\nabla P_1\theta\right)_u  =  \left(\tau_J\nabla
\mathcal{L}_{\mathbb{C}} \theta - i_{\mathbb{C}}\nabla
d_J\theta\right)_u= \left(d_J\mathcal{L}_{\mathbb{C}} \theta -
\mathcal{L}_{\mathbb{C}} d_J\theta\right)_u =
(d_{[J,\mathbb{C}]}\theta)_u=0; \\
\tau_{2}\left(\nabla P_1\theta\right)_u  =  \left(\tau_h\nabla
\mathcal{L}_{\mathbb{C}} \theta - i_{\mathbb{C}}\nabla
d_h\theta\right)_u= \left(d_h\mathcal{L}_{\mathbb{C}} \theta -
\mathcal{L}_{\mathbb{C}} d_h\theta\right)_u =
(d_{[h,\mathbb{C}]}\theta)_u=0; \\
\tau_{3}\left(\nabla P_1\theta\right)_u  =  \left(\tau_J\nabla
d_J\theta\right)_u= \left(d^2_J\theta\right)_u
=\frac{1}{2}(d_{[J,J]}\theta)_u=0; \\
\tau_{4}\left(\nabla P_1\theta\right)_u  =  \left(\tau_h\nabla
d_h\theta\right)_u= \left(d^2_h\theta\right)_u
=\frac{1}{2}(d_{[h,h]}\theta)_u=(d_{R}\theta)_u; \\
\tau_{5}\left(\nabla P_1\theta\right)_u  =  \left(\tau_h\nabla
d_J\theta + \tau_J\nabla d_h\theta\right)_u=
(d_{[h,J]}\theta)_u=0.
\end{gather*}
From the above calculations it follows that a f\/irst-order formal
solution $\theta$ of the system~\eqref{cjh} can be lifted into a
second-order solution if and only if $d_R\theta=0$.
\end{proof}
